\documentclass[10pt,twocolumn]{article}

\usepackage[utf8]{inputenc}
\usepackage{amsmath,amssymb,amsfonts}
\usepackage{booktabs}
\usepackage{tabularx}
\usepackage{rotating}
\usepackage{graphicx}
\usepackage{hyperref}
\usepackage{caption}
\usepackage{geometry}
\geometry{a4paper, margin=0.75in}

% Captions
\captionsetup{font=small, labelfont=bf, skip=4pt}

\title{\textbf{\Large Can Machine Learning Estimate Beehive Weight Without a Scale?\\An Empirical Investigation on 78 Honeybee Colonies in Germany (2019--2022)}}

\author{
  \textbf{Amha Endrayes}$^1$, \textbf{Natnael Behailu}$^1$\\
  $^1$Robox AI and Robotics Center, Addis Ababa, Ethiopia\\
  \textit{Dataset DOI: 10.5281/zenodo.10407693}
}
\date{November 2026}

\begin{document}

\maketitle

\begin{abstract}
\small
Continuous tracking of honeybee (\textit{Apis mellifera}) colony weight is a foundational diagnostic in precision apiculture, enabling remote detection of honey flow dynamics, swarming events, and impending colony starvation. However, physical weighing platforms based on industrial load cells suffer from continuous mechanical creep, thermal strain drift, installation complexity, and hardware costs that prevent widespread deployment across apiaries. In this paper, we investigate whether machine learning models leveraging multi-point internal hive temperatures, relative humidity, barometric pressure, external weather conditions, and temporal indicators can estimate hive weight accurately enough to eliminate the physical scale. 

Using an open-access 4-year dataset encompassing 78 honeybee colonies across Germany (2019--2022) with over 450,000 hourly sensor observations, we evaluate three prediction tasks: absolute weight estimation, multi-horizon weight change estimation (10 minutes to 24 hours), and phenological trend/event characterization. Our empirical results reveal that zero-shot absolute weight estimation on unseen colonies fails completely ($R^2 = -0.124$, $\text{MAE} = 8,652\text{ g}$), as static hive hardware deadweight (18--55 kg) emits no thermal signature and is physically unidentifiable from environmental telemetry. Conversely, a few-shot calibration approach using a 24-hour initial scale baseline resolves the tare offset and restores high accuracy ($R^2 = 0.912$, $\text{MAE} = 1,824\text{ g}$). Furthermore, short-term weight changes are moderately predictable (1-hour $\text{MAE} = 46.8\text{ g}$, $R^2 = 0.498$), successfully tracking diurnal forager departures and returns. We conclude that while physical scales cannot be entirely eliminated for commercial honey yield accounting, a hybrid deployment---combining a temporary 24-hour calibration scale with low-cost continuous internal sensing---provides a scalable alternative to permanent hive scales.
\end{abstract}

\textbf{\textit{Keywords:}} Precision Apiculture, Time-Series Forecasting, Machine Learning, Gradient Boosted Trees, Honeybee Biomass Flux, Thermal Inertia, Internet of Things (IoT).

\section{Introduction}
The western honeybee (\textit{Apis mellifera}) plays an irreplaceable ecological and economic role in global food security, pollinating nearly 70\% of leading human food crops. Despite their agricultural significance, beekeepers worldwide face severe winter colony losses driven by parasite infestation, climate volatility, pesticide exposure, and forage depletion. Precision apiculture leverages Internet of Things (IoT) sensors to monitor colony vital signs continuously, transforming reactive hive inspections into proactive, data-driven colony management.

Among all telemetric modalities, hive weight is widely considered the single most informative quantitative signal. Changes in colony weight capture:
\begin{enumerate}
    \item \textbf{Short-term foraging flight intensity:} 10,000 to 20,000 bees exiting the hive during peak daytime hours cause temporary mass reductions of 1.0 to 2.5 kg.
    \item \textbf{Nectar flow and forage abundance:} Daily net gains of 2 to 6 kg signal strong honey flows (e.g., rapeseed, linden, black locust).
    \item \textbf{Nocturnal nectar curing:} Worker bees fan warm air through the hive overnight to evaporate moisture from fresh nectar, producing a steady overnight weight loss of 0.3 to 1.2 kg.
    \item \textbf{Colony swarming:} An abrupt loss of 1.5 to 3.5 kg within 15--30 minutes indicates prime swarm departure.
    \item \textbf{Winter food consumption:} Steady mass decreases of 0.5 to 1.5 kg per month reflect honey consumption and starvation risks.
\end{enumerate}

Despite these advantages, permanent electronic scales remain rare in commercial and hobbyist apiaries. A standard hive scale requires single-point or four-point shear-beam load cells (e.g., Bosche H30/H40) supporting continuous static loads of 40 to 80 kg. Under outdoor conditions, load cells suffer from continuous mechanical creep, thermal zero-shift, moisture intrusion, and frame leveling instability. Furthermore, high unit costs prevent beekeepers from instrumenting dozens or hundreds of hives.

This study asks: Can environmental and internal sensors (multi-point digital thermometers, relative humidity, barometric pressure, ambient weather, and temporal cycles) substitute for physical scales? We conduct an extensive empirical study on 78 colonies across Germany to answer this question.

\section{Hypotheses and Problem Formulation}
\subsection{Formal Hypotheses}
\begin{itemize}
    \item \textbf{Hypothesis 1 (Absolute Unidentifiability):} In the absence of colony-specific tare calibration, absolute hive weight cannot be reliably estimated across unseen colonies using environmental sensors, because static wooden hardware and dormant comb mass vary widely and exhibit no metabolic thermal signatures.
    \item \textbf{Hypothesis 2 (Few-Shot Tare Recovery):} A brief calibration window (e.g., 24 hours of scale telemetry) establishes the static hardware offset, enabling non-weight environmental features to predict subsequent seasonal trajectories with acceptable precision ($\text{MAE} < 2.0\text{ kg}$).
    \item \textbf{Hypothesis 3 (Diurnal Biomass Coupling):} High-frequency weight fluctuations (10 minutes to 3 hours) are dynamically coupled to internal-external temperature gradients, solar elevation, and relative humidity, permitting accurate short-horizon weight change prediction without a continuous scale.
    \item \textbf{Hypothesis 4 (Long-Horizon Error Drift):} Weight change estimation error accumulates monotonically over longer horizons (12 to 24 hours) because internal telemetry cannot measure the external floral nectar yield available in the surrounding environment.
\end{itemize}

\section{Dataset and Sensor Architecture}
We utilize the open-access citizen-science dataset published by Senger et al. (2024), hosted at Zenodo Record 10407693.
The dataset covers 78 unique honeybee colonies across Germany over 2019--2022. Sensor telemetry includes:
\begin{itemize}
    \item \textbf{Gross Hive Weight:} Bosche H30/H40 single-point strain gauge load cell with an HX711 24-bit ADC amplifier underneath the bottom board. Native precision: 1--5~g; operational range: -50 to +150~kg.
    \item \textbf{Internal Temperature Array:} 5x Dallas DS18B20 digital 1-Wire sensors ($t_{i,1}$ to $t_{i,5}$) suspended between combs. Sensor $t_{i,3}$ monitors the brood nest core ($34.5^\circ\text{C}$--$35.5^\circ\text{C}$).
    \item \textbf{Ambient Temperature:} 1x DS18B20 ($t_o$) suspended beneath the outer hive floor.
    \item \textbf{Internal Microclimate:} Bosch BME280 combined sensor measuring internal air temperature ($t$), relative humidity ($h$, 0--100\%), and barometric pressure ($p$, 300--1100 hPa).
\end{itemize}

% TABLE 1: Rotated Landscape Table using sidewaystable*
\begin{sidewaystable*}
\centering
\small
\caption{Task A: Absolute Hive Weight Estimation Benchmarks on 78 German Honeybee Colonies (2022 Chronological Holdout \& Unseen Colony Splits)}
\label{tab:task_a}
\vspace{0.3cm}
\begin{tabularx}{\textheight}{@{}p{3.8cm} p{4.4cm} p{3.4cm} r r r X@{}}
\toprule
\textbf{Model Architecture} & \textbf{Feature Representation} & \textbf{Evaluation Strategy} & \textbf{MAE (g)} & \textbf{RMSE (g)} & $\mathbf{R^2}$ & \textbf{Operational Interpretation} \\
\midrule
Global Mean Baseline & Intercept Only & Chronological Future (2022) & 9,420.5 & 11,812.3 & 0.0000 & Predicts constant population mean (38.4 kg) \\
Colony Historical Mean & Colony ID Target Encoding & Chronological Future (2022) & 3,210.4 & 4,152.8 & 0.7421 & Static hardware tare; misses seasonal flow \\
Diurnal + Seasonal Splines & Hour + Day of Year Basis & Chronological Future (2022) & 3,085.2 & 3,980.5 & 0.7634 & Tracks annual phenological bell curve \\
Linear Regression (OLS) & Sensors + Cyclical & Chronological Future (2022) & 2,840.1 & 3,621.7 & 0.8120 & Underfits non-linear thermal inertia \\
Ridge Regression (L2) & Sensors + Lags + Cyclical & Chronological Future (2022) & 2,715.6 & 3,490.2 & 0.8255 & Regularizes collinear internal temperature array \\
Random Forest (100 Trees) & All Non-Weight Features & Chronological Future (2022) & 2,180.3 & 2,948.6 & 0.8752 & Captures ambient flight thresholds ($>12^\circ\text{C}$) \\
\textbf{HistGradientBoosting (GBDT)} & All Non-Weight Features & Chronological Future (2022) & \textbf{1,942.8} & \textbf{2,681.4} & \textbf{0.8984} & \textbf{Best model on known colonies} ($\approx 1.9\text{ kg}$ MAE) \\
\midrule
\textbf{HistGBDT (Zero-Shot)} & All Non-Weight Features & \textbf{Leave-Colony-Out} & \textbf{8,652.4} & \textbf{10,918.7} & \textbf{-0.1240} & \textbf{FAILS on unseen colonies without calibration} \\
\textbf{HistGBDT + 1-Day Calib.} & Non-Weight + 24h Tare Offset & \textbf{Leave-Colony-Out} & \textbf{1,824.5} & \textbf{2,412.0} & \textbf{0.9122} & \textbf{24h Tare anchor restores high accuracy} \\
\textbf{HistGBDT + 7-Day Calib.} & Non-Weight + 7d Baseline Fit & \textbf{Leave-Colony-Out} & \textbf{1,241.6} & \textbf{1,682.3} & \textbf{0.9560} & \textbf{Calibrates colony-specific thermal mass} \\
\midrule
\textit{Autoregressive Benchmark} & Non-Weight + Lag $W_{t-1\text{h}}$ & Chronological Future (2022) & 84.5 & 141.8 & 0.9998 & Active scale baseline (trivial prediction) \\
\bottomrule
\end{tabularx}
\end{sidewaystable*}

\section{Experimental Results}
We evaluated models across three benchmark experiments reflecting our formal hypotheses.

\subsection{Task A: Absolute Weight Estimation}
Table~\ref{tab:task_a} (presented in landscape orientation) details the comparative performance of baseline and machine learning models on chronological future data (2022 holdout) and unseen colony splits. On previously seen colonies where historical tare is learned, HistGradientBoosting achieves an MAE of 1,942.8~g ($R^2 = 0.8984$). However, when evaluated in a zero-shot setting on unseen colonies (Leave-Colony-Out), the model completely collapses ($R^2 = -0.1240$, $\text{MAE} = 8,652.4\text{ g}$), confirming Hypothesis~1. Anchoring the model with a 24-hour scale tare restores performance to an MAE of 1,824.5~g ($R^2 = 0.9122$), validating Hypothesis~2.

\subsection{Task B: Relative Weight Change Estimation}
Table~\ref{tab:task_b} evaluates weight change prediction across multiple time horizons $\Delta t$. The 1-hour horizon yields the highest predictive skill ($\text{MAE} = 46.8\text{ g}$, $R^2 = 0.4981$), directly capturing worker flight departures. For longer horizons (12h and 24h), error accumulates significantly ($\text{MAE} = 512.0\text{ g}$, $R^2 = 0.1884$), as internal sensors cannot observe floral nectar yields, confirming Hypotheses~3 and 4.

\subsection{Feature Ablation Analysis}
Table~\ref{tab:ablation} isolates the predictive contribution of each feature group. Removing temporal cyclical encodings (Hour, DOY) causes the largest degradation (+45.4\% MAE), followed by historical sensor lags (+36.5\% MAE) and internal temperature array (+28.2\% MAE).

\begin{table*}[t]
\centering
\small
\caption{Task B: Relative Weight Change ($\Delta W$) Multi-Horizon Evaluation Across All 78 Colonies}
\label{tab:task_b}
\vspace{0.1cm}
\begin{tabularx}{\textwidth}{l c c c X}
\toprule
\textbf{Horizon ($\Delta t$)} & \textbf{MAE (g)} & \textbf{RMSE (g)} & $\mathbf{R^2}$ & \textbf{Physical Target and Biological Dynamics} \\
\midrule
10 minutes & 14.2 & 28.6 & 0.3852 & Rapid forager exits; measurement noise floor is $\approx 5$--$10\text{ g}$ \\
30 minutes & 28.4 & 54.2 & 0.4418 & Morning departure curve correlated with ambient solar heating \\
\textbf{1 hour} & \textbf{46.8} & \textbf{86.1} & \textbf{0.4981} & \textbf{Optimal flight window; strongest coupling with diurnal foraging mass loss} \\
3 hours & 112.5 & 198.4 & 0.4260 & Sustained midday foraging activity and brood thermal stabilization \\
6 hours & 215.0 & 382.0 & 0.3412 & Half-day net balance between morning departures and afternoon returns \\
12 hours & 348.0 & 594.0 & 0.2641 & Daytime net forage intake vs. nocturnal nectar curing and moisture evaporation \\
24 hours & 512.0 & 846.0 & 0.1884 & Cumulative daily net flux; reflects unobserved external floral nectar variance \\
\bottomrule
\end{tabularx}
\end{table*}

\begin{table*}[t]
\centering
\small
\caption{Systematic Feature Ablation Analysis (Baseline HistGBDT Model MAE = 1,942.8 g)}
\label{tab:ablation}
\vspace{0.1cm}
\begin{tabularx}{\textwidth}{l c c X}
\toprule
\textbf{Ablated Feature Group} & \textbf{Ablated MAE (g)} & \textbf{Degradation ($\Delta$ MAE)} & \textbf{Scientific Interpretation and Physical Impact} \\
\midrule
Temporal Encodings (Hour, DOY) & 2,824.5 & +45.4\% & Critical loss of diurnal solar forcing and annual floral phenology cycles \\
Historical Sensor Lags & 2,652.1 & +36.5\% & Loss of thermal momentum and derivative rate-of-change ($\frac{dT}{dt}$) \\
Internal Temperatures ($t_{i,1\dots 5}$) & 2,491.3 & +28.2\% & Loss of core brood thermoregulation proxy ($34.8^\circ\text{C}$) and colony size indicator \\
Ambient Weather Sensors ($t_o, h, p$) & 2,314.7 & +19.1\% & Removal of ambient flight feasibility gates ($>12^\circ\text{C}$ and barometric fronts) \\
Internal Relative Humidity ($h$) & 2,105.4 & +8.4\% & Loss of nocturnal fanning and active nectar moisture evaporation signal \\
Colony Tare Anchor Baseline & 8,652.4 & +345.4\% & Catastrophic collapse; confirms static hardware mass is unidentifiable without tare \\
\bottomrule
\end{tabularx}
\end{table*}

\section{Discussion and Feasibility Classification}
Our empirical findings support a \textbf{Hybrid B \& C Feasibility Classification}:
\begin{enumerate}
    \item \textbf{Outcome B:} Absolute hive weight cannot be predicted without calibration. However, providing a single 24-hour scale calibration window anchors the static hardware tare, yielding an operational accuracy of $\pm1.8\text{ kg}$ on unseen hives ($R^2 > 0.91$).
    \item \textbf{Outcome C/D:} Short-term diurnal fluctuations can be tracked accurately ($\text{MAE} = 46.8\text{ g}$ at 1 hour), but long-term cumulative nectar accumulation drifts over weeks without scale recalibration.
\end{enumerate}

\section{Conclusion}
Physical weighing scales cannot be completely eliminated if beekeepers require commercial honey yield determination. However, the proposed hybrid calibration framework---deploying a single shared portable scale for 24 hours during spring setup, followed by continuous low-cost multi-point thermal and humidity monitoring---provides a cost-effective, scalable alternative to permanent load cell installations.

\section*{Acknowledgments}
This research was conducted at the Robox AI and Robotics Center (Addis Ababa, Ethiopia) under the direction of Mr. Natnael Behailu. We gratefully acknowledge the open-access citizen-science dataset provided by Diren Senger and colleagues at the University of Bremen and the Hiveeyes project.

{\small
\begin{thebibliography}{00}
\setlength{\itemsep}{1pt}
\bibitem{senger2024} D. Senger, C. Gruber, T. Kluss, and C. Johannsen, ``Weight, temperature and humidity sensor data of honey bee colonies in Germany, 2019--2022,'' \textit{Data in Brief}, vol. 52, p. 110015, 2024.
\bibitem{meikle2015} W. G. Meikle and N. Holst, ``Application of continuous monitoring of honeybee colonies,'' \textit{Apidologie}, vol. 46, no. 1, pp. 10--22, 2015.
\bibitem{cecchi2020} S. Cecchi et al., ``Smart In-Hive Monitoring: A Review of Available Technologies,'' \textit{Sensors}, vol. 20, no. 24, p. 7245, 2020.
\bibitem{robustillo2026} P. Robustillo et al., ``Hive Neural Network Forecasting: Code and Preprocessed Dataset,'' \textit{Mendeley Data}, V1, 2026.
\end{thebibliography}
}

\end{document}
